\customchapter{RESUMEN}

Para analizar video largo de forma automática hay que dividirlo en unidades con sentido
propio, y ninguna modalidad debería ocultar a las demás al hacerlo. Esta tesis parte de un
diagnóstico de trabajo previo: en A2Summ-X, una extensión del modelo de resumen A2Summ
\cite{he2023a2summ}, la fusión de modalidades en una representación única hace que el texto
domine la ponderación entre modalidades (pesos de 0.75 a 0.89) y anule la contribución visual
antes de cualquier decisión. Se propone CAVA (\emph{Consumer-Adaptive Video Architecture}),
un sistema de segmentación temática de video largo que no fusiona representaciones. Cada
señal (transcripción, imagen, voz y las descripciones de un modelo visión--lenguaje) emite
sus propias fronteras y etiquetas de tópico. El tópico surge de un consenso de votos con
quórum, las señales sin evidencia en un tramo se abstienen, y cada grupo registra qué
señales lo sostienen.

El sistema produce un mapa de tópicos en dos niveles de granularidad que enlaza segmentos no
adyacentes cuando un tema se retoma. A partir del mapa se generan cinco formatos de versión:
capítulos, recapitulación de 60 segundos, destacados del 10\,\%, hilos de tópico e hilos por
hablante. La validación usa compuertas pre-registradas sobre datos no vistos. En las 47
películas del split de test de MovieNet \cite{huang2020movienet}, el mapa con abstención
supera a la fusión por concatenación ($z = 2.46$) y al gestor contiguo ($z = 2.48$) en MCC
por pares. La ganancia depende de la abstención: sin ella, el consenso no se distingue de la
fusión ($z = 0.80$), y la señal visual sola obtiene un MCC más alto que cualquier
combinación. Las versiones cumplen su presupuesto en todos los tipos de contenido evaluados
(película, reunión, video web y de usuario), y la identidad de hablante alcanza un ARI de 0.73 en
AMI \cite{carletta2007ami}. Tres compuertas no se cumplen: la calidad de los destacados
frente a bases simples (EXP-031/037/038), los enlaces no contiguos juzgados como la misma
trama (G-enlace-2) y la detección del narrador en documentales (EXP-039). De los 26 módulos
del sistema congelado, 10 están validados. La evaluación con personas (EXP-040) está
diseñada, con el conjunto de 9 videos propios ya procesado, y queda pendiente de
pre-registro y ejecución.

\vspace{0.5em}
\noindent \textbf{Palabras clave:}\\
\noindent Segmentación temática de video; agregación multimodal; consenso de votos;
modelos visión--lenguaje; versiones adaptadas al consumidor.
